\section{Background}

The emergence of multi-agent AI systems as operational deployers—rather than mere theoretical constructs—has brought to the fore a set of structural challenges that the field must address before such systems can be relied upon in consequential domains. Among these, the problem of \emph{accountability} is perhaps the most pressing. When a composite of autonomous agents produces a decision or action, the causal chain leading to that outcome is distributed across multiple actors, each contributing in ways that may not be transparent to observers or even to the system designers themselves. This distribution of agency is not a peripheral concern; it strikes at the heart of whether multi-agent systems can be certified for deployment in regulated environments. A regulator reviewing an autonomous system that caused harm needs to be able to trace that outcome to a specific decision node, a specific agent, and ultimately a specific human or institution that authorized or designed that agent's behavior. Without this traceability, the system is accountability-blind by design.

The theoretical foundations for reasoning about accountability in distributed systems trace back to early work on responsibility and agency in multi-agent systems. Dijkstra's 1982 treatment of structured programming and hierarchical system design established that verifiable correctness requires decomposability: a system can only be understood and certified if its behavior can be derived from the behaviors of its constituent parts in a deterministic and inspectable manner \cite{dijkstra1982}. Applied to multi-agent systems, this principle demands that every agent-level decision be traceable to an explicit, auditable state transition—otherwise the system's aggregate behavior remains unverifiable even if individual agents are individually sound. Dijkstra's insight connects directly to what later researchers would characterize as the \emph{tree of accountability}: each leaf node is an agent action, each internal node is a delegation or coordination event, and the root is the final outcome. A break in this chain renders the outcome unresolvable. The challenge in modern multi-agent systems is that this tree is dynamic—it changes shape as agents spawn sub-agents, transfer state, or produce emergent behaviors that were not anticipated in the original design \cite{bansal2019}. Maintaining chain integrity under these conditions is an architectural problem, not merely an organizational one.

Bansal et al. formalize this decomposition in the context of multi-agent systems as a requirement to distinguish between \emph{responsibility}, \emph{blameworthiness}, \emph{accountability}, and \emph{sanctionability}—four conceptually distinct notions that must each be traceable to specific agents or agent groups in order to support regulatory compliance and meaningful oversight \cite{bansal2019}. Responsibility in their framework denotes the duty to perform a given task; blameworthiness attaches after a failure with causal attribution; accountability involves the obligation to explain and justify one's actions to a relevant audience; and sanctionability refers to the consequences that follow from adverse outcomes. Each notion requires different tracing mechanisms and different institutional infrastructure to support it. This decomposition is non-trivial in systems where agents share state, delegate tasks across layer boundaries, or produce emergent behaviors that are not reducible to any single agent's inputs. Without explicit mechanisms to attribute outcomes to specific decision nodes, multi-agent systems remain accountability-blind—a condition that disqualifies them from deployment in healthcare, financial, or infrastructure-critical settings where auditability is a precondition for operation.

Accountability without intervention authority is structurally incomplete. This is where the human-in-the-loop (HITL) paradigm enters the picture. Winner's 1980 treatise on the politics of artifacts established that any automated system operating in a domain affecting human welfare must preserve a designated authority pathway that allows a human supervisor to inspect, override, and correct agent-level decisions in real time \cite{tristr81}. Winner was primarily concerned with the political and institutional implications of autonomous systems—specifically, that automation often shifts power away from human operators in ways that are not visible or reversible. The key design property introduced by this tradition is \emph{reversibility}: a HITL mechanism is only meaningful if the human can actually reverse or halt an ongoing agent action before it produces irreversible consequences. In software engineering terms, this requires that each agent be equipped with a checkpoint protocol that surfaces its current state to the supervisor layer before committing to any downstream effect \cite{tristr81}. Bansal et al. operationalize this further, proposing that HITL integration is not a single design decision but an ongoing process requiring continuous alignment between agent autonomy boundaries and human supervisory capacity as the system's operational context evolves \cite{bansal2019}. This alignment is not static: as agents become more capable and operate at higher speeds, the human supervisory layer faces increasing bandwidth constraints that can degrade the quality of oversight or render it effectively nominal. The design challenge is therefore not merely to include a human in the loop but to ensure that the loop operates at a cadence and granularity sufficient to actually intercept harmful decisions before they propagate.

The fail-safe design tradition extends the HITL principle into the domain of physical and safety-critical systems. The concept originates in control theory and has been codified across multiple engineering disciplines: a system fails safely when its default state upon encountering an undefined condition is a pre-specified safe state rather than undefined or catastrophic behavior. Wiener formalized this principle in 1948 as a core invariant of any feedback-driven system, arguing that the design of automatic controllers must incorporate an explicit boundedness condition guaranteeing that the system either maintains its intended operational envelope or gracefully degrades to a safe mode if that envelope is breached \cite{wiener1948}. This invariant—often referred to as the \emph{fail-safe default}—has since become a foundational requirement in aerospace, nuclear, and medical device standards, and is increasingly treated as a mandatory property for AI-based controllers in those domains. The extension of fail-safe principles to multi-agent AI systems introduces a specific challenge: each agent may implement its own fail-safe logic, but the composite system's emergent behavior upon concurrent agent-level failures is not guaranteed to be safe even if individual agents are individually fail-safe \cite{dijkstra1982}. This is a known limitation that compositional verification discipline helps to address but does not fully resolve, because the interaction surface between agents grows exponentially with the number of agents, making exhaustive verification intractable. What is needed instead is a structural approach that bakes fail-safe behavior into the architecture of agent interaction rather than relying on individual agent implementations.

The BX3 Framework synthesizes these lineages—accountability tracing, HITL intervention authority, and fail-safe defaults—into a three-layer model that structures concern spaces across a consistent abstraction hierarchy. The first layer, \textbf{Intention}, addresses the question of \emph{why} a given agent action was initiated: it captures the goal state, the decision rationale, and the normatively relevant context so that any subsequent audit can reconstruct the purpose behind the action, not merely its mechanical outcome. This layer bears the accountability burden: whatever occupies the Intention layer is answerable for the goals set and the norms invoked. The second layer, \textbf{Execution}, addresses the question of \emph{what} occurred during the transition from intention to outcome: it records state changes, inter-agent communications, and the actual decision points encountered during execution, providing the factual chain of events needed to evaluate whether the transition was faithful to the original intention. This layer is where Dijkstra's compositional traceability requirement lives \cite{dijkstra1982}. The third layer, \textbf{Feedback}, closes the loop by routing outcome signals back to the Intention layer so that goal representations can be updated in light of what actually occurred. Critically, the Feedback layer is where the fail-safe and HITL mechanisms are integrated: anomalous or unexpected outcomes trigger human review or automated revert routines at the Feedback layer before they propagate back into active intention representations. This ensures that the fail-safe property operates at the architectural level rather than being delegated to individual agent implementations \cite{wiener1948}. The three-layer decomposition ensures that accountability, intervention, and fail-safe properties are structurally embedded in how the framework represents and processes agent-level state transitions across all operational stages—making them properties of the architecture itself, not afterthoughts added to an already-designed system.